\input{../doc-class-cours.tex}

\begin{document}

\begin{center}
  \textit{Pourquoi apprendre alors que l’ignorance est instantanée ?} - \textbf{Bill Watterson}
\end{center}

\subsection*{Ex1}
Soit un carré de côté $21cm$. Calculer son périmètre et son aire.

\subsection*{Ex2}
Soit un carré d'aire $140cm^2$. Calculer la longueur de son côté.

\subsection*{Ex3 - Calculer}

\begin{multicols}{3}
\begin{itemize}[label={$\bullet$}]
  \item $25^2 - 20^2$
  \item $100^2$
  \item $(14 - 2)^2$
  \item $\sqrt{81} + \sqrt{36}$
  \item $\sqrt{41} + 9$
  \item $\sqrt{60}$
\end{itemize}
\end{multicols}

\begin{multicols}{2}
\subsection*{PB1}
Les figures sont des carrés. \\
Calculer l'aire $A$.  

\begin{figure}[H]
    \centering
    \includegraphics[width=0.5\linewidth]{4x1-carre/pb1.pdf}
\end{figure}

\subsection*{PB2}
Les figures sont des carrés. \\
Calculer l'aire $A$.  

\begin{figure}[H]
    \centering
    \includegraphics[width=0.5\linewidth]{4x1-carre/pb2.pdf}
\end{figure}
\columnbreak

\subsection*{PB3}
Les figures sont des carrés. \\
On sait que $A_1 = 500cm^2$.\\
Calculer $A_2$.

\begin{figure}[H]
    \centering
    \includegraphics[width=0.5\linewidth]{4x1-carre/pb3.pdf}
\end{figure}

\subsection*{PB4}
Les figures sont des carrés. \\
On sait que $A_1 = 800cm^2$.\\
Calculer AB et BC.

\begin{figure}[H]
    \centering
    \includegraphics[width=0.5\linewidth]{4x1-carre/pb4.pdf}
\end{figure}
\end{multicols}

\end{document}